\chapter*{全局约定}
\addcontentsline{toc}{chapter}{全局约定}

进入具体理论以前，先固定 SI 单位制、时空号规、Fourier 变换和少数跨章反复使用的记号。群论、量子场论、规范理论、散射理论与自由电子量子光学中的专用量在正文首次出现时定义；文末符号索引集中汇总这些约定。

全文采用 SI 单位制，始终显式保留光速 $c$、约化 Planck 常数 $\hbar$、真空介电常数 $\varepsilon_0$、真空磁导率 $\mu_0$ 与 Boltzmann 常数 $k_B$。除人名形成的专有名称外，正文以中文叙述为主；Python 与 TikZ 图内部文字统一使用英文。

\section*{时空与指标}
四维坐标和 Minkowski 度规固定为
\begin{equation*}
 x^\mu=(ct,\mathbf r),
 \qquad
 \eta_{\mu\nu}=\operatorname{diag}(1,-1,-1,-1).
\end{equation*}
因此
\begin{equation*}
 \partial_\mu=(c^{-1}\partial_t,\boldsymbol\nabla),
 \qquad
 \partial^\mu=(c^{-1}\partial_t,-\boldsymbol\nabla).
\end{equation*}
物理四动量统一写为
\begin{equation*}
 p^\mu=(E/c,\mathbf p),
 \qquad
 p^2=\eta_{\mu\nu}p^\mu p^\nu.
\end{equation*}
希腊指标采用 Einstein 求和约定：重复的一上一下指标自动求和；拉丁空间指标取 $1,2,3$。

\section*{Fourier 变换}
时间 Fourier 变换固定为
\begin{equation*}
 \widetilde f(\omega)=\int_{-\infty}^{\infty}\dd t\,f(t)\ee^{+\ii\omega t},
 \qquad
 f(t)=\int_{-\infty}^{\infty}\frac{\dd\omega}{2\pi}\,
 \widetilde f(\omega)\ee^{-\ii\omega t}.
\end{equation*}
一维空间变换采用
\begin{equation*}
 \widetilde f(\kappa)=\int_{-\infty}^{\infty}\dd z\,f(z)\ee^{-\ii\kappa z},
 \qquad
 f(z)=\int_{-\infty}^{\infty}\frac{\dd\kappa}{2\pi}\,
 \widetilde f(\kappa)\ee^{+\ii\kappa z}.
\end{equation*}
波浪号只表示 Fourier 变换。具有逆长度量纲的一般波矢优先使用 $\kappa$ 字母族；真实粒子的物理四动量使用 $p^\mu,k^\mu,q^\mu$ 等动量量纲符号，并在首次出现时分别定义。


高频符号按对象类型固定。通用 Lorentz 因子写作 $\gamma_{\rm L}(\beta)$；中心电子的速度比与 Lorentz 因子固定写作 $\beta_0=v_0/c$ 与 $\gamma_0=\gamma_{\rm L}(\beta_0)$。$\kappa$ 字母族只表示空间波数/波矢；共振线宽与衰减率使用大写 $\Gamma$ 字母族；统计 cumulant 写 $\mathcal C_n$；高斯环境二点函数的对称、反对称核写 $C_{\rm sym}$ 与 $C_{\rm asym}$。无下标的 $\beta$ 只表示经典 PINEM 的无量纲累积耦合；Dirac 矩阵与相对论速度比分别写成 $\beta_D$ 与 $\beta_{\rm rel}$。自由电子量子光学部分的近似误差与小参数统一使用带物理下标的 $\epsilon$ 字母族，探测/收集概率使用 $\mathcal Q^{\rm det}$ 或 $\mathcal Q_{\rm col}$，通道分支比使用 $\mathcal B$。$\eta$ 字母族固定表示 Minkowski 度规：矩阵记作 $\eta$，分量记作 $\eta_{\mu\nu}$；路径积分涨落写成相应场的增量，如 $\delta\varphi$，Dirac Grassmann 外源写作 $\mathcal J_\psi,\overline{\mathcal J}_\psi$。同一物理量在不同平台与近似层级中保持同一记号；近似只删减项或限制参数区间。

\section*{电子电荷与电磁四势}
基本电荷的正值记为 $e>0$，电子电荷始终写
\begin{equation*}
 q_e=-e.
\end{equation*}
电磁四势采用
\begin{equation*}
 A^\mu=(\phi/c,\mathbf A),
 \qquad
 A_\mu=(\phi/c,-\mathbf A).
\end{equation*}
所有公式均保持上述 SI 约定，显式保留 $c$、$\hbar$、$\varepsilon_0$ 与 $\mu_0$；后续章节沿用同一组符号和 Fourier 约定。

\section*{散射与衰变中的单位换算}
涉及粒子实验数值时可以用 eV、MeV 或 GeV 报告能量；它们是能量单位，并不意味着设置 $\hbar=c=1$。物理四动量 $p$ 与能量四矢量 $P_E=cp$ 必须显式区分。若 Mandelstam 不变量 $s$ 用动量平方定义，截面中的面积单位由 $\hbar^2/s$ 提供；若改用 $s_E=c^2s$，同一因子为 $(\hbar c)^2/s_E$。换算时可用
\begin{equation*}
 (\hbar c)^2\simeq0.38938\,\mathrm{mb\,GeV^2},\qquad
 1\,\mathrm{pb}=10^{-40}\,\mathrm{m^2}.
\end{equation*}
衰变率 $\Gamma$ 和能量宽度 $\hbar\Gamma$ 也须区分：前者单位为 $\mathrm{s^{-1}}$，后者单位为能量；单一指数衰变的寿命满足 $\tau=1/\Gamma$。标准模型常数 $G_F^{(E)}$ 具有逆能量平方量纲，SI 拉氏量密度中的系数为 $(\hbar c)^3G_F^{(E)}$。微分截面换变量时必须同时变换测度，例如 $M_{\ell\ell}=c\mathcal Q$ 给 $\dd\mathcal Q^2=2M_{\ell\ell}\dd M_{\ell\ell}/c^2$。
